\chapter*{Preface}
\addcontentsline{toc}{chapter}{Preface}

\textit{Why write this book now?}

The pace of AI development between 2017 and 2026 has been unprecedented. Yet, amidst the flurry of new models and benchmark results, the foundational ideas of the 1940s through 1980s remain surprisingly relevant—and often, underappreciated.

This book emerged from a simple conviction: \textit{you cannot understand where we are going until you understand where we came from}. The transformer did not appear in a vacuum. Its attention mechanism has deep roots in neuroscience and cognitive psychology. The scaling laws of 2020 trace directly back to Rosenblatt's perceptron convergence theorem.

This work is structured as a \textbf{research archive}. Every claim is tied to a primary source. Every equation is re-implemented from scratch in NumPy. Every major paper is analyzed across ten dimensions, from its biological motivation to its modern relevance.

\textit{A note to the reader:} This is not a light history book. It expects familiarity with calculus, linear algebra, and basic probability. However, the mathematical derivations are presented step-by-step, and the scratch code is thoroughly annotated.

I am grateful to the pioneers listed in the biographies section, whose intellectual courage, despite two AI winters, made the modern era possible.

\begin{flushright}
Anshuman Singh\\
India, \today
\end{flushright}